\begin{ExerciseBlock}
\ExerciseHeading{11}{Scheinkräfte auf einem Karussell}{forces_rotating_carousel}
Die horizontale Platte eines Karussells dreht sich in $9\,$s einmal um die Achse (von oben gesehen gegen den Uhrzeigersinn). Ein Schausteller mit der Masse $m=100\,$kg geht zwischen den Sitzen mit einer Geschwindigkeit $v=0,5\,$m/s auf gerader Linie von der Mitte bis zum Rand.

\begin{enumerate}
	\item Berechnen Sie die Winkelgeschwindigkeit des Karussells $\omega$. Geben Sie auch die Richtung von $\vec{\omega}$ an.

	\item Welche Zentrifugalkraft $\vec{F}_{zf}$ verspürt der Schausteller bei $r = 6\,$m? In welche Richtung zieht ihn die Kraft?

	\item Welche Corioliskraft $\vec{F}_{c}$ verspürt der Schausteller? In welche Richtung zieht ihn die Kraft?

	\item Welchen Winkel bildet die Gesamtkraft $\vec{F}_{ges} =  \vec{F}_{zf}  +\vec{F}_{c}$ mit der Radialrichtung? Erstellen Sie eine Skizze.

	\item Skizzieren Sie die Radiusabhängigkeit der Kräfte $F_{zf}(r)$ und $F_{c}(r)$.

\end{enumerate}
\end{ExerciseBlock}
\UseExerciseSolution{forces_rotating_carousel}
